\thispagestyle{empty}
\begin{center}
\vspace*{2.2cm}
{\sffamily\color{azul}\large MANUAL METODOLÓGICO}\\[1.6cm]
{\sffamily\bfseries\fontsize{30}{36}\selectfont\color{tinta}
Jacobianos en el\\[2pt] espacio de secuencias}\\[0.9cm]
{\sffamily\large\color{tinta2}
Cómo resolver, calibrar y estimar modelos macroeconómicos\\[2pt]
y macrofinancieros con agentes heterogéneos}\\[2.4cm]

{\color{gris}\rule{0.55\textwidth}{1.2pt}}\\[1.1cm]

\begin{minipage}{0.80\textwidth}\centering\small\color{tinta2}
Escrito para un estudiante de pregrado en economía que ya sabe optimización
dinámica y programación básica, y que necesita llevar un modelo estático con
heterogeneidad a una versión dinámica que pueda resolverse, calibrarse y
estimarse de verdad.\\[0.7em]
Incluye una implementación propia y verificada del algoritmo \emph{fake news},
y una aplicación completa a un modelo de bancos heterogéneos con encaje,
mercado interbancario y dolarización.
\end{minipage}

\vfill
{\small\color{tinta2}\today}
\end{center}
\clearpage

\thispagestyle{empty}
\section*{Cómo usar este manual}
\addcontentsline{toc}{chapter}{Cómo usar este manual}

Este documento tiene tres niveles de lectura y conviene elegir uno antes de
empezar, en lugar de leerlo de corrido.

\paragraph{Ruta mínima (una tarde)} Capítulo \ref{cap:porque} (por qué existe
el método), sección \ref{sec:jacobiano-def} (qué es exactamente un jacobiano
secuencial), la receta del capítulo \ref{cap:fakenews} y el capítulo
\ref{cap:dag}. Con eso ya se entiende qué hace el código de otra persona.

\paragraph{Ruta de implementación (una o dos semanas)} Todo lo anterior más los
capítulos \ref{cap:bloque}, \ref{cap:extensiones}, \ref{cap:precision} y el
apéndice \ref{ap:codigo}. Al final de esa ruta uno puede escribir su propio
bloque heterogéneo y verificarlo.

\paragraph{Ruta de tesis} Todo, con énfasis en la parte V y el capítulo
\ref{cap:aplicacion}, que traduce paso a paso un modelo estático de bancos
heterogéneos --- con encaje, interbancario y dos monedas --- a un modelo
dinámico resoluble.

\vspace{1em}
A lo largo del texto hay cuatro tipos de recuadro:

\begin{idea}
Lo que hay que recordar aunque se olvide todo lo demás de la sección.
\end{idea}

\begin{trampa}
Errores que cuestan días de depuración y que casi todo el mundo comete al menos
una vez. Vale la pena leerlos aunque uno crea que no los va a cometer.
\end{trampa}

\begin{receta}
Pasos concretos, en orden, para hacer algo. Pensados para seguirse con el
código al lado.
\end{receta}

\begin{tutesis}
Conexiones explícitas con un plan de trabajo sobre heterogeneidad bancaria,
seguro de liquidez y encaje en el Perú. Si ese no es tu tema, estos recuadros
igual sirven como ejemplo de cómo aterrizar el método en un problema propio.
\end{tutesis}

\vspace{1em}
\paragraph{Sobre el código} Todas las figuras de este manual se generaron con
código que acompaña al documento. No hay ningún número inventado: las cifras que
aparecen en el texto salen de ejecutar ese código. La implementación del
algoritmo está escrita desde cero y verificada contra el método directo, que es
lento pero transparente. El apéndice \ref{ap:codigo} explica cómo correrlo.

\paragraph{Sobre la notación} Se sigue de cerca la de
Auclert, Bardóczy, Rognlie y Straub (2021), para que el salto de este manual al artículo original y a su
repositorio sea inmediato. El apéndice \ref{ap:glosario} tiene un glosario
español--inglés de todos los términos técnicos, porque la literatura está en
inglés y a veces la traducción oscurece más de lo que ayuda.
\clearpage
